% !TeX program = xelatex
\documentclass[11pt]{article}
\newcommand{\notesdir}{../../lecture-notes}
\providecommand{\notesdir}{.}

% -----------------------------------------------------------------------------
% Shared preamble for the CO5059 lecture notes and the quiz-review documents.
% Documents outside this folder define \notesdir (path to this folder) first.
% -----------------------------------------------------------------------------

% Page layout and typography
\usepackage[a4paper,portrait,margin=1in]{geometry}
\usepackage{fontspec}
\newfontfamily\arialfont{Arial}
\usepackage{microtype}
\usepackage{parskip}

% Mathematics and diagrams
\usepackage{amsmath,amssymb,amsthm,mathtools}
\usepackage{tikz-cd}
\usetikzlibrary{arrows.meta,positioning,shapes.geometric,shapes.arrows,calc,fit,backgrounds}

% Reusable styles for flowcharts and concept diagrams.
\tikzset{
  concept/.style={
    draw=brandprimary,
    fill=brandsecondary!8,
    rounded corners,
    minimum width=3.2cm,
    minimum height=0.9cm,
    align=center,
    font=\small,
  },
  diagram arrow/.style={-{Latex[length=2.5mm]},thick,brandprimary},
}

% Tables, figures, code, and lists
\usepackage{graphicx}
\graphicspath{{\notesdir/}}
\usepackage{enumitem}
\usepackage{booktabs}
\usepackage{array}
\usepackage{listings}

% Color and links
\usepackage[dvipsnames]{xcolor}
\definecolor{brandprimary}{HTML}{030391}
\definecolor{brandsecondary}{HTML}{1488D8}
\usepackage[
  colorlinks=true,
  linkcolor=brandprimary,
  citecolor=brandsecondary,
  urlcolor=brandsecondary
]{hyperref}

% Headers and footers
\usepackage{fancyhdr}
\setlength{\headheight}{36pt}
\pagestyle{fancy}
\fancyhf{}
\fancyhead[L]{%
  \parbox[c][0.42in][c]{0.42\textwidth}{%
    \raggedright\small\textbf{\coursecode}\\\coursename}}
\fancyhead[R]{%
  \parbox[c][0.42in][c]{0.42\textwidth}{%
    \raggedleft\includegraphics[height=0.42in]{assets/hcmut_logo.png}}}
\fancyfoot[L]{\small Semester \semester}
\fancyfoot[C]{\small\nouppercase{\leftmark}}
\fancyfoot[R]{\small Page \thepage}
\renewcommand{\sectionmark}[1]{\markboth{#1}{}}

% -----------------------------------------------------------------------------
% Course information -- change these values once per course
% -----------------------------------------------------------------------------
\newcommand{\universityname}{HCMC University of Technology}
\newcommand{\faculty}{Faculty of Computer Science and Engineering}
\newcommand{\coursecode}{CO5059}
\newcommand{\coursename}{Advanced System Architectures}
\newcommand{\semester}{261}
\newcommand{\instructor}{Tran Ngoc Thinh}
\newcommand{\instructoremail}{}% leave empty to hide
\newcommand{\notetaker}{Phan Nhat Minh}
\newcommand{\notetakeremail}{minh.phannhat95@hcmut.edu.vn}
\newcommand{\coursedescription}{Quantitative design and evaluation of computer
systems: performance and cost, instruction set architectures, pipelining and
its hazards, memory hierarchy, multiprocessors, superscalar/VLIW processors,
and AI \& security chips.}
\newcommand{\emaillink}[1]{%
  \href{mailto:#1}{\textcolor{brandprimary}{\texttt{#1}}}}

% Colored definitions, results, examples, proofs, and notes.
\input{\notesdir/colored-boxes}

% -----------------------------------------------------------------------------
% Course-specific notation
% -----------------------------------------------------------------------------
\newcommand{\CPI}{\mathrm{CPI}}

\lstdefinestyle{mips}{
  basicstyle=\ttfamily\small,
  keywordstyle=\bfseries\color{brandprimary},
  morekeywords={lw,sw,add,sub,and,or,slt,beq,bne,beqz,xor,nop,addi,addiu,addu,
    subu,li,la,jal,jr,jalr,sll,srl,sra,sllv,srlv,srav,andi,ori,xori,nor,slti,
    sltiu,sltu,lui,mult,multu,div,divu,mfhi,mflo,ll,sc,syscall,move,lh,lhu,lb,
    lbu,sh,sb,blez,bgtz,bltz,bgez,j,mul},
  showstringspaces=false,
  columns=fullflexible,
  keepspaces=true,
}
\lstset{style=mips}

% Creates a consistently formatted lecture heading.
% Usage: \lecture{number}{title}{date}
\newcommand{\lecture}[3]{%
  \section{Lecture #1: #2}
  \label{lec:#1}
  \begin{center}
    \small\textit{#3}
  \end{center}
}


% Cover page + contents page.  Usage: \makecover{Big title}{Description}
\newcommand{\makecover}[2]{%
\pagenumbering{roman}
\begin{titlepage}
  \thispagestyle{empty}

  \begin{center}
    \includegraphics[width=0.82\textwidth]{assets/hcmut_banner_en.png}

    \vspace{0.35cm}
    {\arialfont\large\bfseries\color{brandprimary}\faculty\par}
  \end{center}

  \vfill
  \begin{center}
    {\Huge\bfseries #1\par}
    \vspace{0.55cm}
    {\LARGE\bfseries\color{brandprimary}\coursename\par}
    \vspace{0.2cm}
    {\arialfont\large\bfseries\color{brandsecondary}\coursecode\par}

    \vspace{0.65cm}
    {\color{brandsecondary}\rule{0.38\textwidth}{1.2pt}}
    \vspace{0.65cm}

    {\large\textbf{Instructor:} \instructor
      \ifx\instructoremail\empty\else\ {\small(\emaillink{\instructoremail})}\fi\par}
    \vspace{0.3cm}
    {\large\textbf{Notes by:} \notetaker\
      {\small(\emaillink{\notetakeremail})}\par}
    \vspace{0.3cm}
    {\large\textbf{Semester:} \semester\par}

    \vspace{0.8cm}
    \begin{minipage}{0.72\textwidth}
      \centering\itshape\color{black!70}
      #2
    \end{minipage}
  \end{center}

  \vfill
  \begin{center}
    \small\color{black!60}Last updated: \today
  \end{center}
\end{titlepage}

% -----------------------------------------------------------------------------
% Table of contents -- kept on its own page
% -----------------------------------------------------------------------------
\clearpage
\setcounter{page}{2}
\thispagestyle{empty}
\hypersetup{linkcolor=black}
\setcounter{tocdepth}{2}
\tableofcontents
\hypersetup{linkcolor=brandprimary}
\clearpage
\pagenumbering{arabic}
}


\newcommand{\ex}[3]{\subsection*{Exercise #1 \normalfont\small(book p.~#2, PDF p.~#3)}%
  \addcontentsline{toc}{subsection}{Exercise #1}}

\begin{document}
\sloppy
\makecover{Textbook Exercises: Performance}{Hennessy \& Patterson,
\emph{Computer Architecture: A Quantitative Approach}, 5th ed., Chapter~1
case studies and exercises (1.1--1.18). A map of every exercise with its
page and relevance to Lecture~1, worked solutions for the ones our lectures
cover, and hints for the rest.}

% ============================================================================
\section{How to Use This File}

\begin{itemize}[leftmargin=*]
  \item The exercise text is \textbf{not} copied here -- open the book
  (\texttt{references/Computer Architecture - A Quantitative Approach.pdf})
  at the page given. ``Book p.'' is the printed page number; ``PDF p.'' is
  the page in the PDF viewer.
  \item Each worked solution starts with a one-line summary of what the
  exercise asks, so you can check you are on the right problem.
  \item Chapter~1 of the 5th edition has no plain CPI/clock-cycle
  exercises; those are in Appendix~A (A.1--A.4, see the ISA file) and
  Appendix~C (C.2, C.3, C.7, see the pipelining file).
\end{itemize}

\begin{note}
\textbf{Yield formula.} The 5th edition's text uses
$\text{Yield}=1/(1+D\cdot A)^N$ with $N\approx13.5$ for a 40~nm process, but
that makes the 130/90~nm chips of Exercises 1.1--1.3 have almost zero yield.
The answers below use the classic formula from our lecture,
$\text{Yield}=(1+D\cdot A/4)^{-4}$, which is what these exercises were
written for.
\end{note}

% ============================================================================
\section{Exercise Map}

\begin{center}
\small
\renewcommand{\arraystretch}{1.15}
\begin{tabular}{@{}lll>{\raggedright\arraybackslash}p{5.3cm}l@{}}
\toprule
Ex. & Book p. & PDF p. & Topic & Priority \\
\midrule
1.1 & 62 & 91 & Die yield (Power5); why defect rates differ & \textbf{High} \\
1.2 & 62 & 91 & Dies per wafer, profit per wafer, which chip to make & \textbf{High} \\
1.3 & 62 & 91 & Yield with a spare core; cost per good chip & \textbf{High} \\
1.4 & 63 & 92 & Power supply sizing, disk idle power & Low (power) \\
1.5 & 63--64 & 92--93 & Cooling doors and servers per rack & Low (power) \\
1.6 & 64 & 93 & Discussion: what else matters besides benchmark results & Medium \\
1.7 & 64 & 93 & Dual core at lower frequency/voltage; dynamic power & Low (power) \\
1.8 & 65 & 94 & Moore's law projections, clock-rate limits & Medium \\
1.9 & 65 & 94 & Energy saving for a real-time deadline & Low (power) \\
1.10 & 65--66 & 94--95 & Server-farm power-saving options & Low (power) \\
1.11 & 66 & 95 & FIT, MTTF, availability & Low (not covered) \\
1.12 & 66 & 95 & MTTF of a server farm; cost of downtime & Low (not covered) \\
1.13 & 66 & 95 & SPEC ratios, weighted time ratios, speedup & \textbf{High} \\
1.14 & 66--67 & 95--96 & Amdahl's Law: vectorization & \textbf{High} \\
1.15 & 67 & 96 & Amdahl's Law from a measured fraction of the \emph{new} time & \textbf{High} \\
1.16 & 67--68 & 96--97 & Amdahl's Law with a speedup and a slowdown & \textbf{High} \\
1.17 & 68 & 97 & Amdahl's Law for parallel parts of two applications & \textbf{High} \\
1.18 & 68 & 97 & Amdahl's Law with communication overhead & Medium \\
\bottomrule
\end{tabular}
\end{center}

% ============================================================================
\section{Worked Solutions: High Priority}

\ex{1.1}{62}{91}
\emph{Asks:} the yield of the IBM Power5 from the chip table (Fig.~1.22), and
why its defect rate is lower than the other chips'.
\begin{solution}
\leavevmode
\begin{enumerate}[label=(\alph*),leftmargin=*]
  \item $A = 389\text{ mm}^2 = 3.89\text{ cm}^2$, $D = 0.30/\text{cm}^2$:
  \[
    \text{Yield} = \left(1 + \frac{0.30\times3.89}{4}\right)^{-4} =
    1.292^{-4} \approx 0.36 .
  \]
  \item It is made in an older, more mature 130~nm process. Defect density
  falls as a process matures (the learning curve), while the other two chips
  use a newer 90~nm process.
\end{enumerate}
\end{solution}

\ex{1.2}{62}{91}
\emph{Asks:} profit per 300~mm wafer for a 150~mm$^2$ chip (\$20 profit per
good chip) and a 250~mm$^2$ chip (\$25), both at the Power5's defect rate;
which to build, and how to split 150 wafers/month given the demand.
\begin{solution}
Dies per wafer $= \pi(15)^2/A - \pi\cdot30/\sqrt{2A}$, $D=0.30$.
\begin{enumerate}[label=(\alph*),leftmargin=*]
  \item $A=1.5$: dies $=471.2-54.4 = 416$; yield $=(1+0.1125)^{-4}=0.653$;
  good $\approx 272$; profit $\approx 272\times\$20 \approx \$5{,}430$ per
  wafer.
  \item $A=2.5$: dies $=282.7-42.2 = 240$; yield $=(1+0.1875)^{-4}=0.503$;
  good $\approx 121$; profit $\approx 121\times\$25 \approx \$3{,}020$ per
  wafer.
  \item The smaller chip: it earns $\approx1.8\times$ more per wafer.
  \item Meeting the demand for the small chip needs $50{,}000/272 \approx
  184$ wafers, more than the 150 available, and each of its wafers earns more
  than a large-chip wafer. So all 150 wafers go to the small chip
  ($\approx40{,}700$ chips). The book gives no profit figure for the Power5
  itself, so that sub-question cannot be answered numerically.
\end{enumerate}
\end{solution}

\ex{1.3}{62}{91}
\emph{Asks:} put a second core on the Opteron die and keep a chip if at least
one core is good. Probability of a usable chip, and the new cost per chip if
the old one cost \$20.
\begin{solution}
\leavevmode
\begin{enumerate}[label=(\alph*),leftmargin=*]
  \item One core: $A = 1.99$, $D = 0.75$: $p = (1 + 0.75\times1.99/4)^{-4}
  = 1.373^{-4} \approx 0.28$. At most one of two cores defective:
  \[
    p^2 + 2p(1-p) = 1-(1-p)^2 = 1 - 0.72^2 \approx 0.48 .
  \]
  \item The die doubles in area, so (ignoring edge effects) half as many dies
  fit per wafer and each die costs twice as much before yield. Cost per good
  chip $\approx \$20 \times 2 \times 0.28/0.48 \approx \$23.3$. The spare
  core \emph{raises} the cost: the yield gain (0.28 $\to$ 0.48) does not
  make up for doubling the area.
\end{enumerate}
\end{solution}

\ex{1.13}{66}{95}
\emph{Asks:} choose between Opteron and Itanium~2 for a workload that is 60\%
wupwise-like, 20\% ammp-like, 20\% apsi-like, using the SPEC table
(Fig.~1.17, book p.~43).
\begin{solution}
Data: execution time ratios Opteron/Itanium: wupwise $0.92$, ammp $1.03$,
apsi $0.65$ (times 51.5/56.1, 136/132, 150/231~s).
\begin{enumerate}[label=(\alph*),leftmargin=*]
  \item By overall SPEC (geometric mean of SPECRatios) the Itanium~2 wins:
  27.12 vs.\ 20.86, i.e.\ $1.30\times$.
  \item Weighted average of the time ratios:
  $0.6(0.92) + 0.2(1.03) + 0.2(0.65) = 0.89$.
  \item $1/0.89 \approx 1.13$: for this workload the Opteron is faster,
  the opposite of (a).
\end{enumerate}
\end{solution}
\begin{remark}
A more exact way is to weight the actual times: Opteron $0.6(51.5)+0.2(136)+0.2(150)
= 88.1$~s, Itanium $0.6(56.1)+0.2(132)+0.2(231) = 106.3$~s, a speedup of
$1.21$. Either way the conclusion is the same: the right choice depends on
\emph{your} workload, not on the overall SPEC number.
\end{remark}

\ex{1.14}{66--67}{95--96}
\emph{Asks:} a vector mode $10\times$ faster than normal mode; plot speedup
vs.\ percent vectorization, and find the vectorization needed for various
targets.
\begin{solution}
With vectorized fraction $p$: $\text{Speedup}(p) = 1/\big((1-p) + p/10\big)$.
\begin{enumerate}[label=(\alph*),leftmargin=*]
  \item Speedup vs.\ percent vectorization:
  \begin{center}
  \begin{tikzpicture}[x=6cm,y=0.42cm,font=\scriptsize]
    \draw[->] (0,0) -- (1.08,0) node[right]{Percent vectorization};
    \draw[->] (0,0) -- (0,10.8) node[above]{Net speedup};
    \foreach \x/\l in {0/0,0.2/20,0.4/40,0.6/60,0.8/80,1/100}
      \draw (\x,0) -- (\x,-0.2) node[below]{\l\%};
    \foreach \y in {1,2,4,6,8,10} \draw (0,\y) -- (-0.015,\y) node[left]{\y};
    \draw[thick,brandprimary,domain=0:1,samples=80] plot (\x,{1/((1-\x)+\x/10)});
    \fill[brandsecondary] (0.556,2) circle (1.5pt) node[above left]{(55.6\%, 2)};
    \fill[brandsecondary] (0.889,5) circle (1.5pt) node[above left]{(88.9\%, 5)};
  \end{tikzpicture}
  \end{center}
  The curve stays flat for a long time and shoots up only near 100\%.
  \item $1/2 = 1 - 0.9p \Rightarrow p = 0.5/0.9 = 55.6\%$.
  \item New time $=0.5$; vector part $=p/10 = 0.0556$; share $=
  0.0556/0.5 = 11.1\%$ of the \emph{new} run time.
  \item Half of the maximum (10) is 5: $0.2 = 1 - 0.9p \Rightarrow p = 88.9\%$.
  \item At 70\% with a $20\times$ unit: $1/(0.3 + 0.7/20) = 1/0.335 = 2.99$.
  With the $10\times$ unit the same speedup needs $0.335 = 1 - 0.9p
  \Rightarrow p = 73.9\%$. Raising vectorization by only $\approx4$ points
  matches doubling the vector hardware -- probably the cheaper option.
\end{enumerate}
\end{solution}

\ex{1.15}{67}{96}
\emph{Asks:} an enhancement is $10\times$ faster and is in use for 50\% of
the \emph{enhanced} run time. Overall speedup, and what fraction of the
\emph{original} time was converted.
\begin{solution}
Normalize the new time to 1: 0.5 unenhanced $+$ 0.5 enhanced. The enhanced
half would have taken $0.5\times10 = 5$ originally.
\begin{enumerate}[label=(\alph*),leftmargin=*]
  \item Original time $= 0.5 + 5 = 5.5$, so speedup $= 5.5$.
  \item $5/5.5 = 90.9\%$ of the original time used the fast mode. Check with
  Amdahl: $1/(0.091 + 0.909/10) = 5.5$. \checkmark
\end{enumerate}
\end{solution}
\begin{remark}
This is the classic trap: Amdahl's $F$ is a fraction of the \emph{original}
time. Plugging in $F=0.5$ would give only $1.82$.
\end{remark}

\ex{1.16}{67--68}{96--97}
\emph{Asks:} a faster FP unit ($2\times$, 20\% of time) that slows data-cache
accesses ($1.5\times$ slower, 10\% of time). Overall speedup, and the new
time shares.
\begin{solution}
\leavevmode
\begin{enumerate}[label=(\alph*),leftmargin=*]
  \item $1/(0.8 + 0.2/2) = 1/0.9 = 1.11$.
  \item New time $= 0.7 + 0.2/2 + 0.1\times1.5 = 0.7 + 0.1 + 0.15 = 0.95$;
  speedup $= 1/0.95 = 1.053$.
  \item FP: $0.10/0.95 = 10.5\%$; data cache: $0.15/0.95 = 15.8\%$.
\end{enumerate}
\end{solution}

\ex{1.17}{68}{97}
\emph{Asks:} two applications use 80\% and 20\% of a dual-core's resources;
parallel parts get speedup 2. Speedup of each application alone, and of the
whole system.
\begin{solution}
\leavevmode
\begin{enumerate}[label=(\alph*),leftmargin=*]
  \item App~1, 40\% parallel: $1/(0.6 + 0.4/2) = 1.25$.
  \item App~2, 99\% parallel: $1/(0.01 + 0.99/2) = 1/0.505 = 1.98$.
  \item System, parallelizing app~1: $F = 0.8\times0.4 = 0.32$:
  $1/(0.68+0.16) = 1.19$.
  \item System, parallelizing app~2: $F = 0.2\times0.99 = 0.198$:
  $1/(0.802+0.099) = 1.11$.
\end{enumerate}
Even though app~2 parallelizes almost perfectly, improving app~1 helps the
system more -- it is the common case.
\end{solution}

% ============================================================================
\section{Worked Solutions: Medium Priority}

\ex{1.18}{68}{97}
\emph{Asks:} speedup on $N$ processors for an 80\%-parallel program, then
with communication overhead that grows per processor or per doubling, and the
best $N$.
\begin{solution}
\leavevmode
\begin{enumerate}[label=(\alph*),leftmargin=*]
  \item $1/(0.2 + 0.8/N)$.
  \item Counting 7 added processors: overhead $=7\times0.5\% = 3.5\%$:
  $1/(0.2+0.1+0.035) = 2.99$ (2.94 if you charge all 8 processors).
  \item 1$\to$8 is three doublings, overhead 1.5\%: $1/(0.2+0.1+0.015) =
  3.17$.
  \item $1/\big(0.2 + 0.8/N + 0.005\log_2N\big)$.
  \item Time $T(N) = (1-P) + P/N + 0.005\log_2N$. Setting $dT/dN = -P/N^2 +
  0.005/(N\ln2) = 0$ gives $N^\ast = P/(0.005\ln 2) \approx 289P$ (about 231
  processors for $P = 0.8$).
\end{enumerate}
\end{solution}

\ex{1.8}{65}{94}
\emph{Asks:} project transistor counts and clock rates ten years ahead, and
explain what limited clock growth.
\begin{solution}
(a) Doubling every 2 years for 10 years: $2^5 = 32\times$ (about
$100\times$ with an 18-month doubling). (b)/(c) Apply the 1990s growth rate
and the recent much lower rate from Section~1.4/Fig.~1.11 to the 2005 clock
rate. (d) Power (the power wall): higher clocks need more power than can be
cooled. The extra transistors now go into multiple cores and larger caches.
(e) Apply the stated slowing trend; growth keeps decreasing (roughly halving
again, to the order of 10--20\%/year).
\end{solution}

\ex{1.6}{64}{93}
\emph{Asks:} why someone might choose the server that loses on the reported
power/performance benchmarks.
\begin{solution}
Reported benchmarks may not match your workload (single-thread performance,
floating point, memory size and bandwidth, I/O); also price and
price-performance, software compatibility, reliability/availability,
support, and the fact that a vendor chose which results to report.
\end{solution}

% ============================================================================
\section{Hints: Low Priority (beyond our Lecture~1)}

These use the dynamic power equation and dependability metrics from
Sections~1.5 and 1.7 of the book, which our lecture only mentioned.
\begin{itemize}[leftmargin=*]
  \item \textbf{1.4, 1.5} (pp.~63--64): add the component powers from
  Fig.~1.23, divide by the supply efficiency; servers per cooling door $=$
  door capacity / power per server.
  \item \textbf{1.7, 1.9, 1.10} (pp.~64--66): dynamic power $\propto
  \tfrac12 C V^2 f$; energy $\propto C V^2$. Halving $V$ and $f$ cuts energy
  to $1/4$ and power to $1/8$.
  \item \textbf{1.11} (p.~66): MTTF $=10^9/\text{FIT}$ hours ($10^7$~h for
  FIT 100); availability $=\text{MTTF}/(\text{MTTF}+\text{MTTR})$; with 1000
  parts that all fail together, MTTF$/1000$.
  \item \textbf{1.12} (p.~66): treat failures as independent with rate
  $1/\text{MTTF}$ per machine; compare the extra hardware cost with the
  downtime cost from Fig.~1.3.
\end{itemize}

\end{document}
